\section{Data, Calendar, and Occupational Harmonization}\label{sec:data}

\subsection{CPS stocks and the corrected calendar}

The public-data application uses IPUMS Basic Monthly CPS records from January 2017 through July 2026. Employed respondents with positive final weights are aggregated into occupation-by-age-group-by-month stocks. Ages 22--25 form the young group and ages 26--65 the comparison group. Age is not labor-market experience, and the outcome can change through employment entry, exit, occupational switching, or the older denominator.

The original wide extract accidentally selected the March ASEC sample (`03s') rather than the Basic Monthly sample (`03b') in 2017--2021. ASEC records carry zero Basic Monthly final weights for this purpose. A separately authenticated extract restores those five March samples. October 2025 remains absent because no CPS was collected during the federal shutdown. The combined files contain 9,843,021 records and 114 observed survey months. Omitting December 2022 as a transition leaves 113 months in the substantive static model.

The corrected calendar is now the substantive baseline. Restoring March changes the coefficient from $-0.1311$ to $-0.1346$, with a 9,999-draw interval $[-0.2223,-0.0468]$. The frozen result remains visible as the confirmatory chronology record but no longer defines the default exposition.

\subsection{Taxonomy bridge and universe reconciliation}

CPS occupations use Census-2010 codes through 2019 and Census-2018 codes thereafter, while source exposures use SOC vintages. Pre-2020 records are expanded along official Census-2010-to-2018 routes and aggregated to the common target taxonomy. In a one-to-many route, official conversion proportions are applied equally to young and older source records. This preserves the source occupation's age ratio across target components by construction; official files do not supply age-specific probabilities.

The full route-expanded preperiod universe contains 539 target occupations with positive young and older stock. Two counts that appeared to be a five-occupation discrepancy are instead nonnested, architecture-specific supports. The stored 490-occupation universe requires full beta coverage. The 495-occupation continuous-mapping universe requires repaired AIOE and Webb coverage. They overlap in 466 occupations: 24 appear only in beta support and 29 only in AIOE-plus-Webb support. Adding Webb to beta leaves the 468-occupation primary model. The online appendix names every member and reason.

Within primary support, split sources contribute 14.93 percent of young and 17.75 percent of older weighted stock during 2017--2019. Ninety-nine target occupations can receive a split route and carry 23.38 percent of fitted target information; proportional allocation of information by split-stock share assigns 4.44 percent to split stock. The former is a conservative touched-target diagnostic and the latter is not an identified decomposition.

The negative association does not depend on split routes. Restricting the corrected model to 369 targets with no structurally possible split inbound route gives $-0.1674$ under fixed labels and $-0.1471$ after recomputing the comparison on that support. Aggregating to a stable Census-2010 taxonomy gives $-0.1522$. Each changes the population and should not be called a correction of the primary estimate.

Age-specific route allocation remains uncertain. Under unrestricted official-route accounting, the 2017--2019 young/older stock ratio can range from .1031 to .1480 in Q1 and from .0790 to .1175 in Q5. These are sharp accounting ranges for stocks, not bounds on the nonlinear regression coefficient. A predeclared tilt that changes young-versus-older allocation odds toward the highest-beta eligible target by factors from .5 to 2 moves the coefficient from $-0.1292$ to $-0.1391$. Only 6.88 percent of early stock is eligible for that scenario, so it is not a global robustness claim.

\subsection{A detectable workflow failure}

Literal exact-code matching of SOC-2010 AIOE into a post-2018 outcome taxonomy covers just 3.33 percent of computer-and-mathematical employment. The documented bridge reaches 97.7 percent. A machine-readable compatibility gate now rejects a merge when declared source and target vintages are not identical and no versioned route is supplied.

The corresponding continuous sequence separates layers: original score/original 410-occupation support gives $-0.01885$; repaired score/same support gives $-0.01920$; repaired score/expanded 495-occupation support gives $-0.03156$; excluding computer/math gives $-0.02940$. The invalid exact-code result is not an economically defensible alternative. The small first movement concerns score correction; the larger second movement is population re-admission.

\subsection{Treatment formation}

The primary Eloundou beta score and Webb software control are standardized using preperiod employment weights. Beta is divided into employment-weighted quintiles with ties retained. On the 468-occupation primary support, Q1 contains 133 occupations and 20.02 percent of employment; Q5 contains 91 and 19.53 percent. All support and sensitivity rows state whether cutoffs and scales are frozen or recomputed. Respondent-equivalent counts after route expansion are constructed records, not unique people.

\begin{table}[!htbp]
\centering
\caption{Reconciliation of occupation universes}\label{tab:sampleflow}
\begin{threeparttable}
\small
\begin{tabularx}{\textwidth}{@{}YrrY@{}}
\toprule
Universe & Occupations & Months used to form support & Availability rule \\
\midrule
Raw route-expanded balanced candidates & 539 & 66 & Positive young and older preperiod stock \\
Frozen beta-complete support & 490 & 66 & Raw balance plus finite beta \\
Repaired AIOE plus Webb support & 495 & 66 & Raw balance plus finite repaired AIOE and Webb \\
Intersection of preceding two supports & 466 & 66 & Satisfies both architecture-specific rules \\
Primary beta plus Webb model & 468 & 108/113 & Finite beta and Webb in frozen/corrected static panel \\
Six core measures plus Webb & 444 & 108 & Literal complete-case intersection \\
\bottomrule
\end{tabularx}
\tabnote{The 490- and 495-occupation supports are not nested: 24 occupations occur only in beta-complete support and 29 only in repaired-AIOE-plus-Webb support. The first four rows define preperiod availability; the final two are fitted-model samples after the full calendar and model-specific requirements. Frozen and corrected calendars contain 108 and 113 static-model months, respectively.}
\end{threeparttable}
\end{table}

